% Options for packages loaded elsewhere
\PassOptionsToPackage{unicode}{hyperref}
\PassOptionsToPackage{hyphens}{url}
\PassOptionsToPackage{dvipsnames,svgnames,x11names}{xcolor}
\documentclass[
  10pt,
  fontset=fandol]{ctexart}
\usepackage{xcolor}
\usepackage[margin=16mm]{geometry}
\usepackage{amsmath,amssymb}
\setcounter{secnumdepth}{-\maxdimen} % remove section numbering
\usepackage{iftex}
\ifPDFTeX
  \usepackage[T1]{fontenc}
  \usepackage[utf8]{inputenc}
  \usepackage{textcomp} % provide euro and other symbols
\else % if luatex or xetex
  \usepackage{unicode-math} % this also loads fontspec
  \defaultfontfeatures{Scale=MatchLowercase}
  \defaultfontfeatures[\rmfamily]{Ligatures=TeX,Scale=1}
\fi
\usepackage{lmodern}
\ifPDFTeX\else
  % xetex/luatex font selection
\fi
% Use upquote if available, for straight quotes in verbatim environments
\IfFileExists{upquote.sty}{\usepackage{upquote}}{}
\IfFileExists{microtype.sty}{% use microtype if available
  \usepackage[]{microtype}
  \UseMicrotypeSet[protrusion]{basicmath} % disable protrusion for tt fonts
}{}
\makeatletter
\@ifundefined{KOMAClassName}{% if non-KOMA class
  \IfFileExists{parskip.sty}{%
    \usepackage{parskip}
  }{% else
    \setlength{\parindent}{0pt}
    \setlength{\parskip}{6pt plus 2pt minus 1pt}}
}{% if KOMA class
  \KOMAoptions{parskip=half}}
\makeatother
\usepackage{longtable,booktabs,array}
\newcounter{none} % for unnumbered tables
\usepackage{calc} % for calculating minipage widths
% Correct order of tables after \paragraph or \subparagraph
\usepackage{etoolbox}
\makeatletter
\patchcmd\longtable{\par}{\if@noskipsec\mbox{}\fi\par}{}{}
\makeatother
% Allow footnotes in longtable head/foot
\IfFileExists{footnotehyper.sty}{\usepackage{footnotehyper}}{\usepackage{footnote}}
\makesavenoteenv{longtable}
\setlength{\emergencystretch}{3em} % prevent overfull lines
\providecommand{\tightlist}{%
  \setlength{\itemsep}{0pt}\setlength{\parskip}{0pt}}
\usepackage{amsmath,amssymb,mathtools,needspace}
\xeCJKDeclareCharClass{CJK}{"2032,"0400->"04FF,"2160->"216B,"2460->"2473,"25A0,"25C7,"25CB,"25CF}
\IfFontExistsTF{Arial Unicode MS}{\xeCJKsetup{AutoFallBack=true}\setCJKfallbackfamilyfont{\CJKrmdefault}{Arial Unicode MS}}{}
\setcounter{secnumdepth}{0}
\setlength{\emergencystretch}{2em}
\allowdisplaybreaks
\usepackage{bookmark}
\IfFileExists{xurl.sty}{\usepackage{xurl}}{} % add URL line breaks if available
\urlstyle{same}
\hypersetup{
  pdftitle={习题},
  colorlinks=true,
  linkcolor={Maroon},
  filecolor={Maroon},
  citecolor={Blue},
  urlcolor={Blue},
  pdfcreator={LaTeX via pandoc}}

\title{习题}
\author{}
\date{}

\begin{document}
\maketitle

\subsection{习题}\label{ux4e60ux9898}

\textbf{1.} 确定下列求积公式中的待定参数，使其代数精度尽量高，并指明所构造出的求积公式所具有的代数精度：

（1）

\[
\int_{-h}^{h}f(x)\,\mathrm{d}x\approx A_{-1}f(-h)+A_0f(0)+A_1f(h);
\]

（2）

\[
\int_{-2h}^{2h}f(x)\,\mathrm{d}x\approx A_{-1}f(-h)+A_0f(0)+A_1f(h);
\]

（3）

\[
\int_{-1}^{1}f(x)\,\mathrm{d}x\approx[f(-1)+2f(x_1)+3f(x_2)]/3;
\]

（4）

\[
\int_0^h f(x)\,\mathrm{d}x\approx h[f(0)+f(h)]/1+ah^2[f'(0)-f'(h)].
\]

\textbf{2.} 分别用梯形公式和 Simpson 公式计算下列积分：

（1）

\[
\int_0^1\frac{x}{4+x^2}\,\mathrm{d}x,\quad n=8;
\]

（2）

\[
\int_0^1\frac{(1-e^{-x})^{\frac12}}{x}\,\mathrm{d}x,\quad n=10;
\]

（第 2 题其余小题续下页。）

\begin{quote}
校注（agent 补充，CH04B-E005）：第 1 题（4）首项的分母在原扫描中明确印为``1''，故保留 \(h[f(0)+f(h)]/1\)。该处疑为原书排印错误：仅作排印一致性检查，当 \(f\equiv1\) 时，积分为 \(h\)，而所印右端为 \(2h\)，并且导数项为零；疑应将分母印为``2''。此处不求解题中参数或代数精度。\href{../assets/ch04b/review-p117-exercises.png}{习题原扫描局部}。
\end{quote}

\href{../../source-images/pdf-118.jpeg}{原页图像}

\subsection{习题（续）}\label{ux4e60ux9898ux7eed}

\textbf{2.（续）}

（3）

\[
\int_1^9\sqrt{x}\,\mathrm{d}x,\quad n=4;
\]

（4）

\[
\int_0^{\pi/6}\sqrt{4-\sin^2\varphi}\,\mathrm{d}\varphi,\quad n=6.
\]

\textbf{3.} 直接验证 Cotes 公式（4.2.4）具有五次代数精度。

\textbf{4.} 用 Simpson 公式求积分 \(\displaystyle\int_0^1 e^{-x}\,\mathrm{d}x\) 并估计误差。

\textbf{5.} 推导下列三种矩形求积公式：

\[
\int_a^b f(x)\,\mathrm{d}x=(b-a)f(a)+\frac{f'(\eta)}2(b-a)^2,
\]

\[
\int_a^b f(x)\,\mathrm{d}x=(b-a)f(b)-\frac{f'(\eta)}2(b-a)^2,
\]

\[
\int_a^b f(x)\,\mathrm{d}x=(b-a)f\left(\frac{a+b}2\right)+\frac{f''(\eta)}{24}(b-a)^3.
\]

\textbf{6.} 证明梯形公式（4.2.9）和 Simpson 公式（4.2.11）当 \(n\to\infty\) 时收敛到积分 \(\displaystyle\int_a^b f(x)\,\mathrm{d}x\)。

\textbf{7.} 用复化梯形公式求积分 \(\displaystyle\int_a^b f(x)\,\mathrm{d}x\)。要将积分区间 \([a,b]\) 分成多少等份，才能保证误差不超过 \(\varepsilon\)（不计舍入误差）？

\textbf{8.} 用 Romberg 方法计算积分 \(\displaystyle\frac2{\sqrt\pi}\int_0^1 e^{-x}\,\mathrm{d}x\)，要求误差不超过 \(10^{-5}\)。

\textbf{9.} 卫星轨道是一个椭圆，椭圆周长的计算公式是

\[
s=a\int_0^{\pi/2}\sqrt{1-\left(\frac ca\right)^2\sin^2\theta}\,\mathrm{d}\theta,
\]

其中 \(a\) 是椭圆的长半轴，\(c\) 是地球中心与轨道中心（椭圆中心）的距离。记 \(h\) 为近地点距离，\(H\) 为远地点距离，\(R=6\,371\ \mathrm{km}\) 为地球半径，则 \(a=(2R+H+h)/2,\quad c=(H-h)/2\)。

我国第一颗人造卫星近地点距离 \(h=439\ \mathrm{km}\)，远地点距离 \(2\,384\ \mathrm{km}\)，试求卫星轨道的周长。

\textbf{10.} 证明等式

\[
n\sin\frac\pi n=\pi-\frac{\pi^3}{3!\,n^2}+\frac{\pi^5}{5!\,n^4}-\cdots,
\]

试依据 \(n\sin\dfrac\pi n\)（\(n=3,6,12\)）的值，用外推算法求 \(\pi\) 的近似值。

\textbf{11.} 用下列方法计算积分 \(\displaystyle\int_1^3\frac{\mathrm{d}y}{y}\)，并比较结果。

（1）Romberg 方法；

（2）三点及五点 Gauss 公式；

（3）将积分区间分为四等份，用复化两点 Gauss 公式。

\textbf{12.} 用三点公式和五点公式求 \(f(x)=\dfrac1{(1+x)^2}\) 在 \(x=1.0,1.1\) 和 \(1.2\) 处的导数值，并估计误差。\(f(x)\) 的值由表 4.9 给出。

\textbf{表 4.9}

{\def\LTcaptype{none} % do not increment counter
\begin{longtable}[]{@{}
  >{\centering\arraybackslash}p{(\linewidth - 10\tabcolsep) * \real{0.1667}}
  >{\centering\arraybackslash}p{(\linewidth - 10\tabcolsep) * \real{0.1667}}
  >{\centering\arraybackslash}p{(\linewidth - 10\tabcolsep) * \real{0.1667}}
  >{\centering\arraybackslash}p{(\linewidth - 10\tabcolsep) * \real{0.1667}}
  >{\centering\arraybackslash}p{(\linewidth - 10\tabcolsep) * \real{0.1667}}
  >{\centering\arraybackslash}p{(\linewidth - 10\tabcolsep) * \real{0.1667}}@{}}
\toprule\noalign{}
\begin{minipage}[b]{\linewidth}\centering
\(x\)
\end{minipage} & \begin{minipage}[b]{\linewidth}\centering
\(1.0\)
\end{minipage} & \begin{minipage}[b]{\linewidth}\centering
\(1.1\)
\end{minipage} & \begin{minipage}[b]{\linewidth}\centering
\(1.2\)
\end{minipage} & \begin{minipage}[b]{\linewidth}\centering
\(1.3\)
\end{minipage} & \begin{minipage}[b]{\linewidth}\centering
\(1.4\)
\end{minipage} \\
\midrule\noalign{}
\endhead
\bottomrule\noalign{}
\endlastfoot
\(f(x)\) & \(0.250\,0\) & \(0.226\,8\) & \(0.206\,6\) & \(0.189\,0\) & \(0.173\,6\) \\
\end{longtable}
}

\begin{quote}
校注（agent 补充，CH04B-E006）：第 9 题的原扫描明确写 \(s=a\int_0^{\pi/2}\cdots\,\mathrm{d}\theta\)，本转录保留原式。若 \(s\) 表示文中的椭圆``周长''，该式疑漏系数 \(4\)：所印表达式对应四分之一周长；以 \(c=0\) 的圆作一致性检查，原式给出 \(\pi a/2\)，而圆周长为 \(2\pi a\)。仅记录原书公式疑误，不代做本题数值计算。\href{../assets/ch04b/review-p118-orbit.png}{局部原扫描}。
\end{quote}

\subsection{本节覆盖原页的校注记录}\label{ux672cux8282ux8986ux76d6ux539fux9875ux7684ux6821ux6ce8ux8bb0ux5f55}

下列记录按原页关联，含同页相邻小节的校注；计算时同时读取上文校注和完整原页。

\begin{itemize}
\tightlist
\item
  \texttt{CH04B-E005}：原书 PDF 页 117
\item
  \texttt{CH04B-E006}：原书 PDF 页 118
\end{itemize}

\end{document}
